\chapter{社会效益与长期发展战略}

\section{项目社会效益}

\begin{enumerate}[leftmargin=1.6em]
  \item \hlgreen{助力粮食安全}：通过精准识别与窗口期预警，减少病虫害导致的减产与绝收，保障主粮稳产增收。
  \item \hlgreen{推动农药减量}：以「对症下药、精准用药」替代「盲目用药、过量用药」，从源头减少农药滥用与残留，促进绿色农业与食品安全。
  \item \hlgreen{弥合数字鸿沟}：把专业植保能力以极简工具送到亿万小农户手中，让「数字农业」真正惠及最需要的人群，助力乡村振兴与共同富裕。
  \item \hlgreen{赋能基层农技}：为基层农技人员提供数字化工具，缓解「人少面广」的困境，提升农技服务的效率与覆盖面。
\end{enumerate}

\section{长期发展愿景}

本项目的长期愿景是成为\hlgold{「中国小农户的第一农事参谋」}：以病虫害诊断这一高频刚需为入口，逐步延伸至水肥管理、农事提醒、产销对接等更多农事场景，构建覆盖「种—管—防—收」全流程的轻量化数字农事服务体系，让每一个小农户都能用一部手机，获得专业的农事决策支持。

\section{业务拓展规划}

\begin{enumerate}[leftmargin=1.6em]
  \item \hlgreen{横向拓场景}：由病虫害诊断延伸至缺素诊断、水肥管理、生育期识别等农事场景，形成全流程农事参谋。
  \item \hlgreen{纵向拓作物}：由主粮作物向高价值经济作物（果蔬、中药材、茶叶）拓展，提升单用户价值。
  \item \hlgreen{数据拓价值}：沉淀区域病虫害发生地图与农资需求数据，向农资供应链、农业保险、政府决策提供数据服务，开辟第二增长曲线。
  \item \hlgreen{区域拓市场}：由国内县域样板向全国农业主产区，乃至「一带一路」沿线农业国家输出产品与能力。
\end{enumerate}

本项目相信，技术的力量在于普惠。以「轻量化的 AI」服务「最广大的小农户」，既是一条清晰的商业路径，更是一项有意义的社会事业。
